%!TEX program = xelatex
\documentclass[cn,hazy,blue,14pt,normal]{elegantnote}
\title{笔记模板}

\author{···}
\institute{···}

\date{\zhtoday}

\usepackage{array}

\begin{document}

\maketitle

\newpage

\begin{itemize}
  \item 适配不同设备，包括 Pad（默认），Screen（幻灯片），Kindle，PC（双页），通用（A4 纸张）；
  \item 5 套颜色主题，分别是：\textcolor{eblue}{blue}（默认）、\textcolor{egreen}{green}、\textcolor{ecyan}{cyan}、 \textcolor{sakura}{sakura} 和 \textcolor{black}{black}；
  \item 语言支持：中文（默认），英文；
  \item 全局字体大小支持：8pt, 9pt, 10pt, 11pt, 12pt, 14pt, 17pt, 和 20pt；
\end{itemize}

\begin{remark}
  如果你想为自己的文档添加底色，可以在导言区添加下面设置：
  \begin{lstlisting}
    \definecolor{geyecolor}{RGB}{199,237,204}
    \pagecolor{geyecolor}
  \end{lstlisting}
\end{remark}

\subsection{设备选择}

为了让笔记方便在不同设备上阅读，免去切边，缩放等操作，本模板适配不同的设备，分别为 Pad（默认），Kindle，PC，A4。
不同屏幕的选择为
\begin{lstlisting}[frame=none]  
  \documentclass[device=pad]{elegantnote}    % ipad screen size
  \documentclass[device=kindle]{elegantnote} % kindle screen size
  \documentclass[device=pc]{elegantnote}     % double pages for pc 
  \documentclass[device=normal]{elegantnote} % a4 normal page
  \documentclass[device=screen]{elegantnote} % 4:3 PPT size
\end{lstlisting}

\subsection[颜色主题]{颜色主题\footnote{测试章节脚注。}}

本模板内置 5 套颜色主题，分别是 \textcolor{eblue}{blue}（默认），\textcolor{egreen}{green}， \textcolor{ecyan}{cyan}， \textcolor{sakura}{sakura} 和 \textcolor{black}{black}。如果不需要颜色，可以选择黑色（black）主题。颜色主题的设置方法：
\begin{lstlisting}[frame=none]  
  \documentclass[green]{elegantnote}
  \documentclass[color=green]{elegantnote}
  ...
  \documentclass[black]{elegantnote}
  \documentclass[color=black]{elegantnote}
\end{lstlisting}

\begin{note}
只有中文模式才可输入中文，如果需要在英文模式下输入中文，可以自行添加 \lstinline{ctex} 宏包\footnote{需要使用 \lstinline{scheme=plain} 选项才不会把标题改为中文。}或者使用 \lstinline{xeCJK} 宏包设置字体。另外如果在笔记中使用了抄录环境（\lstinline{lstlisting}），并且里面有中文字符，请务必使用 \hologo{XeLaTeX} 编译。
\end{note}


\subsection{定理类环境}

此模板采用了 \lstinline{amsthm} 中的定理样式，使用了 4 类定理样式，所包含的环境分别为
\begin{itemize}
  \item \textbf{定理类}：theorem，lemma，proposition，corollary；
  \item \textbf{定义类}：definition，conjecture，example；
  \item \textbf{备注类}：remark，note，case；
  \item \textbf{证明类}：proof。
\end{itemize}

\begin{remark}
在选用 \lstinline{lang=cn} 时，定理类环境的引导词全部会改为中文。
\end{remark}


\section{写作示例}

我们将通过三个步骤定义可测函数的积分。首先定义非负简单函数的积分。以下设 $E$ 是 $\mathcal{R}^n$ 中的可测集。

\begin{definition}[可积性]
设 $ f(x)=\sum\limits_{i=1}^{k} a_i \chi_{A_i}(x)$ 是 $E$ 上的非负简单函数，其中 $\{A_1,A_2,\ldots$, $A_k\}$ 是 $E$ 上的一个可测分割，$a_1,a_2,\ldots,a_k$ 是非负实数。定义 $f$ 在 $E$ 上的积分为 1. 3
\begin{equation}
   \label{inter}
   \int_{E} f dx = \sum_{i=1}^k a_i m(A_i).
\end{equation}
一般情况下 $0 \leq \int_{E} f dx \leq \infty$。若 $\int_{E} f dx < \infty$，则称 $f$ 在 $E$ 上可积。
\end{definition}

一个自然的问题是，Lebesgue 积分与我们所熟悉的 Riemann 积分有什么联系和区别？之后我们将详细讨论 Riemann 积分与 Lebesgue 积分的关系。这里只看一个简单的例子。设 $D(x)$ 是区间 $[0,1]$ 上的 Dirichlet 函数。即 $D(x)=\chi_{Q_0}(x)$，其中 $Q_0$ 表示 $[0,1]$ 中的有理数的全体。根据非负简单函数积分的定义，$D(x)$ 在 $[0,1]$ 上的 Lebesgue 积分为
\begin{equation}\label{inter2}
  \int_0^1 D(x)dx = \int_0^1 \chi_{Q_0} (x) dx = m(Q_0) = 0
\end{equation}
即 $D(x)$ 在 $[0,1]$ 上是 Lebesgue 可积的并且积分值为零。但 $D(x)$ 在 $[0,1]$ 上不是 Riemann 可积的。

\begin{table}[htbp]
  \centering
  \small
  \caption{燃油效率与汽车价格}
    \begin{tabular}{lcc}
    \toprule
                  &       (1)         &        (2)      \\
    \midrule
    燃油效率      &   -238.90***      &      -49.51     \\
                  &    (53.08)        &      (86.16)    \\
    汽车重量      &                   &        1.75***  \\
                  &                   &       (0.641)   \\
    常数项        &  11253.00***      &    1946.00      \\
                  &  (1171.00)        &   (3597.00)     \\
    观测数        &     74            &      74         \\
    $R^2$         &      0.220        &       0.293     \\
    \bottomrule
    \end{tabular}%
  \label{tab:reg}%
\end{table}%

\begin{theorem}[Fubini 定理]\label{thm:fubi}
若 $f(x,y)$ 是 $\mathcal{R}^p\times\mathcal{R}^q$ 上的非负可测函数，则对几乎处处的 $x\in \mathcal{R}^p$，$f(x,y)$ 作为 $y$ 的函数是 $\mathcal{R}^q$ 上的非负可测函数，$g(x)=\int_{\mathcal{R}^q}f(x,y) dy$ 是 $\mathcal{R}^p$ 上的非负可测函数。并且
\begin{equation}\label{eq:461}
  \int_{\mathcal{R}^p\times\mathcal{R}^q} f(x,y) dxdy=\int_{\mathcal{R}^p}\left(\int_{\mathcal{R}^q}f(x,y)dy\right)dx.
\end{equation}
\end{theorem}

\begin{proof}
Let $z$ be some element of $xH \cap yH$.  Then $z = xa$ for some $a \in H$, and $z = yb$ for some $b \in H$. If $h$ is any element of $H$ then $ah \in H$ and $a^{-1}h \in H$, since $H$ is a subgroup of $G$. But $zh = x(ah)$ and $xh = z(a^{-1}h)$ for all $h \in H$. Therefore $zH \subset xH$ and $xH \subset zH$, and thus $xH = zH$.  Similarly $yH = zH$, and thus $xH = yH$, as required.
\end{proof}

\end{document}
